% !TeX root = ../main.tex
% Kapiteldatei: Erfahrungen

\chapter{Erfahrungen bei Aufbau, Verkabelung, Inbetriebnahme und Steuerung der CNC-Heißdrahtschneidanlage}

\section*{Hinweis zur Erstellung dieses Erfahrungsberichts}

Während des gesamten Projekts wurden KI-Werkzeuge unterstützend eingesetzt. Die KI diente dabei nicht nur zur Recherche, sondern auch zur Dokumentation von Arbeitsschritten, Problemen, Lösungsansätzen und Zwischenergebnissen während der Entwicklung der Maschine.

Der vorliegende Erfahrungsbericht basiert daher auf den tatsächlich während des Projekts entstandenen Notizen, Diskussionen, Messergebnissen und dokumentierten Problemen. Die textliche Ausformulierung und Strukturierung wurde durch KI unterstützt. Sämtliche beschriebenen Erfahrungen, Entscheidungen, Schwierigkeiten und technischen Erkenntnisse stammen jedoch aus der praktischen Projektarbeit und wurden von den Teammitgliedern überprüft, ergänzt und freigegeben.

Die KI wurde somit als Werkzeug zur Aufbereitung und Zusammenfassung der Projekterfahrungen verwendet, nicht zur Erzeugung der zugrunde liegenden Inhalte.

\section{Erfahrungen im Projektverlauf}

Im Verlauf des Projekts wurde deutlich, dass der Aufbau einer CNC-Heißdrahtschneidanlage nicht nur aus der mechanischen Konstruktion besteht, sondern in hohem Maße von einer zuverlässigen elektrischen Verdrahtung, einer systematischen Inbetriebnahme und einem guten Verständnis der Steuerungssoftware abhängt. Besonders im Bereich der Elektronik, der GRBL-Konfiguration und der praktischen Fehlersuche war der Lernaufwand deutlich größer als zunächst erwartet.

Zu Beginn lag der Fokus vor allem darauf, die grundlegenden Komponenten miteinander in Betrieb zu nehmen. Verwendet wurden ein Arduino Uno R3, ein CNC Shield V3, Schrittmotortreiber, NEMA17-Schrittmotoren, ein 24-V-Netzteil, Endschalter, ein PWM-MOSFET-Modul sowie später der Heizdraht. Obwohl diese Komponenten auf den ersten Blick gut dokumentiert und weit verbreitet wirken, zeigte sich in der Praxis, dass gerade bei günstigen CNC-Shield-Varianten, Klon-Platinen und unterschiedlichen Treibermodulen viele Details beachtet werden müssen. Die Beschriftung auf den Platinen, die tatsächlich verwendeten Pins, die GRBL-Version und die Pinbelegung des CNC Shields stimmten nicht immer intuitiv mit den Erwartungen überein.

Ein wesentlicher erster Schritt war die Inbetriebnahme der Schrittmotoren. Anfangs lag trotz angeschlossener 24-V-Versorgung kein Haltemoment an den Motoren an. Der entscheidende Punkt war dabei die Enable-Brücke zwischen EN und GND auf dem CNC Shield. Erst nachdem diese Brücke gesetzt wurde, wurden die Treiber aktiv und die Motoren erhielten Haltemoment. Diese Erfahrung zeigte früh, dass selbst kleine Jumper oder scheinbar nebensächliche Steckbrücken über die komplette Funktionsfähigkeit des Systems entscheiden können.

Die ersten Motortests wurden zunächst mit einfachen Arduino-Testsketchen durchgeführt. Dabei konnten die STEP- und DIR-Pins der X-, Y- und Z-Achse erfolgreich angesteuert werden. Später wurde GRBL 1.1h installiert und mit Universal Gcode Sender verbunden. Auch hier gab es zunächst Schwierigkeiten, da zuerst eine ungeeignete bzw. veraltete GRBL-Bibliothek verwendet wurde. Erst nach der Installation der offiziellen GRBL-Version und dem Hochladen des Beispiels grblUpload meldete sich der Arduino korrekt mit Grbl 1.1h ['\$' for help]. Dieser Schritt war entscheidend, weil dadurch die Grundlage für eine standardisierte CNC-Steuerung geschaffen wurde.

Während der Inbetriebnahme zeigte sich außerdem eine grundlegende Einschränkung der gewählten Steuerungsplattform, die in der frühen Projektphase nicht ausreichend berücksichtigt worden war. Ursprünglich war vorgesehen, die Maschine mit vier vollständig unabhängigen Achsen zu betreiben. Diese Annahme beruhte unter anderem auf Vorbildern ähnlicher Heißdrahtschneidanlagen, die häufig einen Arduino Mega 2560 in Verbindung mit einem RAMPS-Board verwenden. In solchen Konfigurationen kommen häufig spezielle GRBL-Varianten oder andere Firmwarelösungen zum Einsatz, welche tatsächlich vier unabhängig interpolierbare Achsen unterstützen.

Im Projekt wurde jedoch ein Arduino Uno R3 mit CNC Shield V3 eingesetzt. Zu Beginn wurde davon ausgegangen, dass die vier Treibersteckplätze des Shields auch vier unabhängig steuerbare Achsen ermöglichen würden. Erst während der intensiveren Beschäftigung mit GRBL und der tatsächlichen Pinbelegung wurde deutlich, dass das auf dem Arduino Uno eingesetzte Standard-GRBL 1.1h lediglich drei Bewegungsachsen unterstützt. Der vierte Treibersteckplatz (A-Achse) dient standardmäßig lediglich als Klon einer vorhandenen Achse und kann nicht ohne Weiteres als eigenständige vierte CNC-Achse betrieben werden.

Rückblickend handelte es sich hierbei weniger um ein technisches Problem als vielmehr um eine Fehleinschätzung in der Planungsphase. Innerhalb des Teams wurde die Funktionalität verschiedener Hardwareplattformen teilweise gleichgesetzt, obwohl sich Arduino Uno mit CNC Shield V3 und Arduino Mega 2560 mit RAMPS sowohl hinsichtlich der verfügbaren Ein- und Ausgänge als auch hinsichtlich der eingesetzten Firmware erheblich unterscheiden. Die Unterschiede wurden erst im Verlauf der praktischen Inbetriebnahme vollständig verstanden.

Für das konkrete Projekt stellte sich diese Einschränkung letztlich als akzeptabel heraus, da die rechte vertikale Achse als Klon der linken vertikalen Achse betrieben werden konnte. Dennoch war diese Erkenntnis ein wichtiger Lernschritt. Sie verdeutlichte, wie entscheidend eine frühzeitige Überprüfung der tatsächlichen Fähigkeiten einer Steuerungsplattform ist und dass äußerlich ähnliche Systeme nicht zwangsläufig dieselben Funktionen bereitstellen. Daraus ergab sich die finale Achslogik: X und Z bilden die beiden horizontalen Achsen der linken und rechten Seite, Y und A bilden die beiden vertikalen Achsen, wobei A als Klon von Y arbeitet.

Trotz der beschriebenen Einschränkungen der eingesetzten Steuerungsplattform wurde die Maschine von Beginn an mit vier mechanisch unabhängigen Achsen konstruiert. Jede Achse verfügt über einen eigenen Antrieb und kann grundsätzlich separat verfahren werden. Die Beschränkung auf drei unabhängig interpolierbare Achsen ergibt sich daher nicht aus der mechanischen Konstruktion, sondern ausschließlich aus der Kombination aus Arduino Uno R3, CNC Shield V3 und der darauf eingesetzten Standardversion von GRBL.

\section{Erfahrungen im Projektverlauf}

Im Verlauf des Projekts wurde bewusst entschieden, die vorhandene Steuerungsplattform beizubehalten, da sie für die grundlegende Funktion der Heißdrahtschneidanlage ausreichend war und eine zuverlässige Inbetriebnahme ermöglichte. Gleichzeitig wurde bei der Konstruktion darauf geachtet, zukünftige Erweiterungen nicht auszuschließen.

Die Maschine kann daher als hardwareseitig für eine spätere Vierachssteuerung vorbereitet betrachtet werden. Durch den Austausch der Steuerung gegen leistungsfähigere Plattformen, beispielsweise auf Basis eines Arduino Mega 2560 mit RAMPS-Board oder anderer moderner CNC-Steuerungen, wäre grundsätzlich eine Ansteuerung aller vier Achsen als vollständig unabhängige Bewegungsachsen möglich. Die hierfür erforderlichen Anpassungen betreffen überwiegend die Elektronik und Firmware, während die mechanische Konstruktion weitgehend unverändert weiterverwendet werden könnte.

Rückblickend erwies sich dieser Umstand sogar als Vorteil. Obwohl die ursprünglich angestrebte Vierachssteuerung mit der gewählten Hardware nicht realisiert werden konnte, entstand eine mechanische Plattform, die über die tatsächlich genutzte Funktionalität hinaus erweiterbar bleibt. Dadurch bietet die Maschine nicht nur die Möglichkeit zum Heißdrahtschneiden in der aktuell verwendeten Konfiguration, sondern stellt gleichzeitig eine geeignete Basis für zukünftige Entwicklungen und weiterführende Untersuchungen im Bereich mehrachsiger CNC-Systeme dar.

Ein großer Teil der Arbeit bestand aus Fehlersuche. Besonders prägend war der Ausfall der Motorversorgung auf einem CNC Shield. Obwohl am Eingang des Shields stabile 24 V gemessen wurden, hatten die Motoren plötzlich kein Haltemoment mehr, die Treiber blieben kalt und GRBL nahm Befehle zwar an, bewegte aber nur intern die Positionswerte. Durch systematisches Messen per Multimeter wurde festgestellt, dass am VMOT-Pin der Treiber keine 24 V mehr anlagen. Schließlich wurde ein kleines glasartiges Sicherungsbauteil auf dem Shield als unterbrochen identifiziert. Eine provisorische Überbrückung führte sofort wieder zu Haltemoment, wodurch die Fehlerursache eindeutig auf die unterbrochene VMOT-Versorgung eingegrenzt werden konnte. Diese Fehlersuche war sehr lehrreich, weil sie zeigte, dass Software, Arduino und GRBL völlig korrekt arbeiten können, während die Leistungsversorgung der Motoren trotzdem unterbrochen ist.

Auch ein späterer Kurzschluss machte deutlich, wie empfindlich solche Aufbauten sein können. Besonders kritisch sind das Umstecken von Motoren oder Treibern bei eingeschalteter 24-V-Versorgung, lose Brücken im Leistungsstrompfad und ungesicherte Messungen mit Messspitzen. Daraus ergab sich als feste Arbeitsregel: Vor jedem Umstecken oder Ändern an Motoren, Treibern, Endschaltern, Jumpern oder Leistungsverdrahtung müssen 24 V ausgeschaltet, USB getrennt und einige Sekunden gewartet werden. Diese Regel wurde im weiteren Verlauf konsequent beachtet.

Ein weiterer wichtiger Punkt war die korrekte Zuordnung der Motorleitungen. Bei Schrittmotoren ist nicht entscheidend, welche Leitung als „Plus“ oder „Minus“ bezeichnet wird, sondern dass die beiden Spulenpaare korrekt zusammen auf dem jeweiligen Treiberausgang liegen. Mit dem Multimeter wurden die Leitungen auf Durchgang geprüft, um die beiden Spulenpaare zu identifizieren. Ein falsches Mischen der Spulenleitungen kann dazu führen, dass ein Motor nur ruckelt, gar nicht läuft oder im schlimmsten Fall Treiber belastet werden. Die Drehrichtung kann dagegen einfach durch Drehen eines Motorsteckers bzw. Vertauschen eines Spulenpaars oder alternativ über GRBL-Parameter angepasst werden.

Besonders zeitintensiv war die Kalibrierung der Achsen. Zunächst war unklar, wie Microstepping, mechanische Übersetzung, Riemensteigung, Spindelsteigung und GRBL-Parameter zusammenspielen. Entscheidend wurde schließlich die Erkenntnis, dass GRBL die Microstepping-Jumper auf dem Shield nicht automatisch erkennt. Die Hardware-Jumper verändern nur das Verhalten des Treibers; softwareseitig muss dies über die Schritte pro Millimeter angepasst werden. Die relevanten Werte sind \$100 für X, \$101 für Y und \$102 für Z. Durch praktische Messfahrten wurde die Maschine schließlich auf sinnvolle Werte gebracht. Für die aktuelle Konfiguration mit 1/2-Mikroschrittauflösung ergaben sich als Kalibrierwerte X = 10 Schritte/mm, Y = 50 Schritte/mm und Z = 10 Schritte/mm. A übernimmt als Y-Klon entsprechend den Y-Wert.

Bei der Kalibrierung kam es zwischenzeitlich zu großer Verwirrung, weil nach Änderungen der Steps/mm-Werte teilweise die jeweils erste Fahrbewegung nicht dem erwarteten Verhalten entsprach. Dadurch wurden anfangs Messwerte aufgenommen, die scheinbar nichtlinear wirkten und zu widersprüchlichen Schlussfolgerungen führten. Erst durch wiederholtes, systematisches Messen mit definierten Jog-Bewegungen wurde klar, welche Werte tatsächlich reproduzierbar waren. Diese Erfahrung war wichtig, weil sie zeigt, dass bei CNC-Inbetriebnahmen einzelne Messungen oft nicht ausreichen. Man muss Bewegungen mehrfach testen, reproduzierbar messen und nach jeder Änderung prüfen, ob die Maschine wirklich im erwarteten Zustand ist.

Auch die mechanische Laufruhe war ein großes Thema. Im Vollschrittbetrieb und bei niedrigen Vorschüben traten deutliche Schwingungen und Resonanzen auf. Teilweise wurde der gesamte Rahmen angeregt. Zunächst lag der Verdacht nahe, dass die Motoren oder Treiber falsch eingestellt waren. Später zeigte sich jedoch, dass die Schwingungen stark von der Drehzahl, der Beschleunigung, dem Microstepping und der mechanischen Reibung (Stick-Slip-Effekte) abhingen. Besonders auffällig war, dass sehr langsame Bewegungen teilweise deutlich schlechter klangen als höhere Vorschübe. Bei höheren Feed Rates liefen die Motoren wesentlich ruhiger, während bei bestimmten Zwischenbereichen wieder Resonanzen auftraten. Dies entspricht dem typischen Verhalten von Schrittmotoren, die in bestimmten Drehzahlbereichen mechanische Resonanzen anregen können.

Eine große Verbesserung brachte das Schmieren der Spindeln. Diese waren zuvor vollständig entfettet und liefen trocken. Nach dem Auftragen von Langzeitfett liefen die Spindeln deutlich ruhiger. Wahrscheinlich wirkte das Fett nicht nur reibungsmindernd, sondern auch dämpfend zwischen Spindel und Mutter. Dadurch wurde die Geräuschentwicklung deutlich reduziert. Diese Erfahrung zeigte, dass elektrische oder softwareseitige Probleme nicht isoliert betrachtet werden dürfen. Mechanik, Elektronik und Steuerung beeinflussen sich gegenseitig.

Im Bereich der GRBL-Einstellungen wurden zahlreiche Parameter angepasst und getestet. Besonders wichtig waren \$1=255, damit die Motoren dauerhaft bestromt bleiben und Haltemoment behalten, die Steps/mm-Werte \$100, \$101, \$102, die Maximalgeschwindigkeiten \$110, \$111, \$112 und die Beschleunigungen \$120, \$121, \$122. Die Beschleunigung wurde von einem sehr konservativen Wert auf 30 mm/s² erhöht, um schneller aus stark ratternden Niedrigdrehzahlbereichen herauszukommen. Die maximale Geschwindigkeit wurde auf 2000 mm/min gesetzt. Diese Werte erwiesen sich für die Maschine als sinnvoller Startpunkt.

Die Endschalter und das Homing waren ebenfalls nicht trivial. Die Endschaltermodule wurden über ihre Kontakte geprüft. Je nach Anschluss und Logik zeigte GRBL über Pn:X, Pn:Y oder Pn:XYZ an, welche Endschalter als aktiv erkannt wurden. Dabei wurde deutlich, dass \$5 für die Invertierung der Limit-Pins entscheidend ist. Außerdem zeigte sich bei CNC Shield V3 und GRBL 1.1h eine mögliche Besonderheit im Zusammenhang mit Z-Limit und Spindle/PWM-Pins. In Foren wurde beschrieben, dass bei manchen Shield-Versionen die Zuordnung zwischen Z-Endschalter und Spindle-Pin durch Änderungen in GRBL 1.1 problematisch sein kann. Diese Beobachtung passte teilweise zu den eigenen Problemen mit dem Z-Endschalter. Deshalb wurde gelernt, nicht blind auf die Beschriftung eines Shields zu vertrauen, sondern im Zweifel mit Multimeter und GRBL-Statusmeldungen zu prüfen, welcher Pin tatsächlich welche Funktion hat.

Die Steuerung über Universal Gcode Sender war insgesamt sehr hilfreich, brachte aber ebenfalls Lernaufwand mit sich. UGS wurde genutzt, um GRBL zu verbinden, Einstellungen auszulesen, Achsen zu joggen und G-Code-Programme auszuführen. Dabei wurde klar, dass UGS und die Arduino IDE nicht gleichzeitig denselben COM-Port nutzen können. Wenn keine Verbindung möglich war, lag dies häufig an einem blockierten Port, einem offenen seriellen Monitor, einem falschen COM-Port oder einem Arduino, auf dem GRBL nicht korrekt lief. Bei Upload-Problemen half es, das CNC Shield vom Arduino zu entfernen und den Arduino separat zu testen. Dadurch konnte zwischen Arduino-/Firmware-Problemen und Shield-/Hardwareproblemen unterschieden werden.

Das Arbeiten mit G-Code wurde Schritt für Schritt aufgebaut. Zunächst wurden nur einzelne Bewegungen ausgeführt, später Quadrate, Diagonalen und komplexere Testformen. Dabei wurde gelernt, dass jede G-Code-Zeile einen eigenen Bewegungsauftrag darstellt. Wenn mehrere Achsen in derselben Zeile stehen, interpoliert GRBL diese Bewegung so, dass die Achsen gleichzeitig am Ziel ankommen. Stehen Bewegungen in getrennten Zeilen, werden sie nacheinander ausgeführt. Für den Heißdrahtschneider ist dies entscheidend, weil X und Z bewusst gleiche oder unterschiedliche Wege fahren können. Dadurch lässt sich der Draht horizontal schräg stellen, was später für komplexere Schnitte und unterschiedliche Profile relevant ist.

Auch die Bedeutung grundlegender G-Code-Befehle wurde im Verlauf des Projekts praktisch erarbeitet. Der Befehl G21 stellt die Maßeinheit auf Millimeter ein, G90 aktiviert absolute Koordinaten und G94 legt fest, dass Vorschübe in Millimetern pro Minute angegeben werden. Für relative Bewegungen wurde G91 verwendet. Der Befehl M3 S... aktiviert den Heißdraht mit einer vorgegebenen Leistung, während M5 den Heißdraht wieder abschaltet. Mit G4 P... kann eine definierte Wartezeit eingefügt werden, beispielsweise um dem Heizdraht Zeit zum Aufheizen zu geben. Außerdem zeigte sich, dass der Vorschubwert F möglichst in jedem Programm explizit gesetzt werden sollte, damit keine unbeabsichtigt übernommenen Vorschubwerte aus vorherigen Programmen verwendet werden.

Die Integration des Heißdrahtsystems stellte einen weiteren wichtigen Schritt dar. Der Heizdraht wird nicht direkt vom Arduino angesteuert, sondern über ein PWM-fähiges MOSFET-Modul geschaltet. Dabei war zunächst unklar, welcher Ausgang des CNC Shields tatsächlich das von GRBL erzeugte PWM-Signal bereitstellt. Erst nach intensiver Recherche und mehreren Tests konnte die korrekte Zuordnung verifiziert werden. Wichtig ist dabei, dass das PWM-Signal des Arduino mit dem Signaleingang des MOSFET-Moduls verbunden wird und dass Arduino und MOSFET eine gemeinsame Masseverbindung (GND) besitzen.

Für die Leistungssteuerung des Heizdrahts wurden mehrere GRBL-Parameter verwendet. Der Parameter \$30=1000 definiert den maximal möglichen Leistungswert. Der Parameter \$31=0 definiert den minimalen Leistungswert. Dadurch ergibt sich ein Leistungsbereich von 0 bis 1000. Wird beispielsweise der Befehl M3 S1000 ausgeführt, entspricht dies 100 \% der eingestellten Ausgangsleistung. M3 S500 entspricht 50 \% und M3 S0 beziehungsweise M5 entspricht 0 \% Leistung. Auf diese Weise kann die Heizleistung des Drahtes über den S-Wert im G-Code feinstufig angepasst werden.

Zusätzlich wurde der Parameter \$32 betrachtet. Mit \$32=1 wird der sogenannte Laser Mode von GRBL aktiviert. Obwohl diese Funktion ursprünglich für Laseranwendungen entwickelt wurde, eignet sie sich auch für Heißdrahtschneider, da der Heizdraht ebenfalls eine kontinuierliche Energiequelle darstellt. In diesem Modus werden Bewegungen flüssiger geplant und die Leistungsregelung während der Bewegung besser an die tatsächliche Verfahrgeschwindigkeit angepasst. Dadurch können gleichmäßigere Schneidbedingungen erreicht werden als im normalen Betriebsmodus (\$32=0).

Beim Heizdraht zeigte sich, dass schon kleine PWM-Werte zu hoher Temperatur führen können. Bei einer Drahtlänge von etwa 600 mm und 0,4 mm Durchmesser, Werkstoff NiCr, kann bei 24 V sehr viel Leistung entstehen. Deshalb wurden Werte im Bereich von etwa S50 bis S150 als sinnvoller Startbereich betrachtet. Wenn der Draht bereits bei sehr kleinen S-Werten extrem heiß wird, muss geprüft werden, ob das MOSFET-Modul tatsächlich per PWM regelt oder dauerhaft durchschaltet. Der Befehl M5 muss den Draht zuverlässig abschalten.

Ein OLED-Display mit SSD1306-Controller wurde ebenfalls betrachtet. Es stellte sich jedoch heraus, dass Standard-GRBL auf dem Arduino Uno keine direkte Displayunterstützung bietet. Der Speicher des Uno ist bereits durch GRBL stark ausgelastet. Eine Integration eines Displays in GRBL wäre daher aufwendig und für den Projektabschluss nicht notwendig. Sinnvoller wäre ein separater Mikrocontroller für Anzeigezwecke oder später eine leistungsfähigere Steuerung wie grblHAL oder FluidNC. Für die aktuelle Maschine bleibt UGS am Laptop die praktikablere Lösung.

Insgesamt zeigte das Projekt sehr deutlich, dass die Inbetriebnahme einer selbstgebauten CNC-Maschine einen erheblichen praktischen Lernanteil besitzt. Viele Probleme lassen sich nicht allein durch Datenblätter oder Tutorials lösen, sondern müssen durch Messen, Testen, Ausschließen und systematisches Vorgehen eingegrenzt werden. Besonders wichtig war dabei, immer nur eine Änderung auf einmal vorzunehmen und danach das Verhalten zu prüfen. Sobald mehrere Parameter gleichzeitig geändert wurden, wurde die Fehlersuche deutlich schwieriger.

Rückblickend waren die größten Herausforderungen nicht einzelne Bauteile, sondern das Zusammenspiel aller Teilsysteme: Mechanik, Motoren, Treiber, Shield, Arduino, GRBL, UGS, Endschalter, PWM-MOSFET und Heizdraht. Erst als verstanden wurde, welche Aufgabe jedes Teilsystem übernimmt und wie es mit den anderen zusammenhängt, konnte die Maschine zuverlässig betrieben werden.

Ein besonders prägendes Erlebnis war die abschließende Vorführung der Maschine. Nach einer intensiven Arbeitsphase in den letzten Stunden vor der Präsentation konnte die Anlage schließlich vollständig montiert, verkabelt und in Betrieb genommen werden. Sowohl die Achsbewegungen als auch die Heißdrahtsteuerung funktionierten einwandfrei, und es konnten einfache wie auch komplexere Formen aus Styropor mittels G-Code geschnitten werden. Kurz vor der Vorführung im Rahmen der Abgabe der Hardware traten jedoch erneut Probleme auf. Obwohl das System wenige Stunden zuvor noch funktionsfähig gewesen war, kam es zu weiteren Hardwareausfällen und Kontaktproblemen. Unter anderem musste kurzfristig ein Arduino ersetzt werden, zusätzlich traten Fehler in der Verkabelung sowie am PWM-Modul für den Heizdraht auf.

Dadurch konnte die Maschine während der offiziellen Vorführung nicht in dem Zustand präsentiert werden, in dem sie zuvor bereits erfolgreich betrieben worden war. Dies war zunächst enttäuschend, verdeutlichte jedoch gleichzeitig eine wichtige Erkenntnis des Projekts: Gerade bei komplexen Prototypen entscheidet nicht nur die grundsätzliche Funktion, sondern auch die langfristige Zuverlässigkeit und Robustheit der Gesamtkonstruktion. Die dabei gewonnenen Erfahrungen hinsichtlich Fehlersuche, Verkabelung, Dokumentation und systematischer Inbetriebnahme waren letztlich ebenso wertvoll wie die erfolgreiche Realisierung der Maschine selbst.

Das Projekt erforderte daher nicht nur mechanisches Konstruieren und Montieren, sondern auch Kenntnisse in Elektrotechnik, Fehlersuche, Firmware, CNC-Steuerung, G-Code und praktischer Inbetriebnahme. Besonders wertvoll war, dass viele Fehler real auftraten und gelöst werden mussten: fehlendes Haltemoment, falsche Treiberfreigabe, defekte bzw. unterbrochene Motorversorgung, falsch interpretierte Endschalterzustände, unklare Shield-Pinbelegung, Upload-Probleme, COM-Port-Probleme, Resonanzen, falsch kalibrierte Achsen, heiße Treiber und PWM-Probleme am Heizdraht.

Gerade diese Schwierigkeiten haben das Verständnis für CNC-Systeme stark vertieft. Am Ende war die Maschine nicht nur mechanisch aufgebaut, sondern auch steuerungstechnisch verstanden. Die Achsen konnten kalibriert, Bewegungen interpoliert, G-Code-Programme ausgeführt, Homing vorbereitet und der Heizdraht über PWM angesteuert werden. Damit wurde aus einzelnen Komponenten eine funktionierende CNC-Heißdrahtschneidanlage.

Für zukünftige Projekte würde aus diesen Erfahrungen folgen, von Anfang an noch konsequenter modular zu testen: zuerst den reinen Arduino mit GRBL, dann das Shield ohne Treiber, dann einzelne Treiber und Motoren, danach Endschalter, danach alle Achsen, danach PWM und erst zuletzt der Heizdraht. Außerdem sollten Schaltplan, Pinbelegung, Jumperstellungen und GRBL-Parameter während der Inbetriebnahme fortlaufend dokumentiert werden. Gerade bei mehreren Umbauten, Ersatzteilen und Fehlersuchen ist eine solche Dokumentation entscheidend, um später nachvollziehen zu können, welcher Zustand tatsächlich funktioniert hat.

Zusammenfassend war die Inbetriebnahme anspruchsvoller und zeitintensiver als erwartet, aber auch der lehrreichste Teil des Projekts. Die aufgetretenen Probleme waren nicht nur Hindernisse, sondern führten zu einem wesentlich besseren Verständnis der Maschine. Besonders der Umgang mit GRBL, UGS, Achskalibrierung, Endschaltern, Treibern und PWM-Heißdrahtsteuerung war für das Projekt zentral. Die gewonnenen Erfahrungen stellen daher einen wesentlichen Bestandteil des Projektergebnisses dar. Neben der fertigen Maschine liefert insbesondere die Dokumentation der aufgetretenen Probleme, Fehlannahmen und Lösungswege wertvolle Erkenntnisse für zukünftige Projekte.
